\documentclass[conference]{IEEEtran}

% Vietnamese text and typography
\usepackage[utf8]{inputenc}
\usepackage[T5]{fontenc}
\usepackage{lmodern}

% Mathematics, tables, and links
\usepackage{amsmath,amssymb}
\usepackage{booktabs}
\usepackage{multirow}
\usepackage{array}
\usepackage{graphicx}
\usepackage{float}
\usepackage{xcolor}
\usepackage[hidelinks]{hyperref}

\newcommand{\meanstd}{trung bình $\pm$ độ lệch chuẩn}
\newcommand{\best}[1]{\textbf{#1}}

\title{Báo cáo Thực nghiệm Bổ sung cho Mô hình\\
MoE--MobileNetV3 trên Tập dữ liệu PlantDoc}

\author{%
\IEEEauthorblockN{Van-Hau Bui, Duc-Thang Doan, Duc-Minh Tran, and Trong-Minh Hoang}
\IEEEauthorblockA{Posts and Telecommunications Institute of Technology, Hanoi, Vietnam\\
University of Economics--Technology for Industries (UNETI), Hanoi, Vietnam\\
Faculty of Data Science and Artificial Intelligence, College of Technology, National Economics University, Hanoi, Vietnam\\
Corresponding author: Trong-Minh Hoang (e-mail: hoangtrongminh@ptit.edu.vn)}}

\begin{document}
\maketitle

\begin{abstract}
Báo cáo này tổng hợp tám nhóm thực nghiệm bổ sung phục vụ quá trình
revision của mô hình Mixture-of-Experts (MoE) dựa trên MobileNetV3-Small.
Các thí nghiệm lần lượt kiểm tra tính nhất quán của kiến trúc, mức sử dụng
expert, vai trò của routing phụ thuộc đầu vào, chi phí thực thi sparse động,
ảnh hưởng của Top-$k$, đối chứng static-uniform có dung lượng tương đương,
ảnh hưởng của context 6 chiều và baseline MobileNetV3-Small đồng nhất về
implementation. Ablation có kiểm soát cho thấy Top-2 đạt Accuracy
$0.8596\pm0.0124$ và Macro-F1 $0.8275\pm0.0153$, cao nhất trong các cấu hình
$k\in\{1,2,3,4\}$. Nhánh context 6-D tạo chênh lệch trung bình $0.0267$
Accuracy và $0.0257$ Macro-F1 so với zero-context, đồng thời có độ lệch
chuẩn thấp hơn trên cả hai metric.
\end{abstract}

\begin{IEEEkeywords}
Mixture-of-Experts, MobileNetV3, dynamic sparse routing, Top-$k$,
PlantDoc, ablation study, context-aware routing.
\end{IEEEkeywords}

\section{Mục tiêu và phạm vi}

Mục tiêu của đợt thực nghiệm là cung cấp bằng chứng định lượng cho các câu
hỏi Q1, Q4, Q7 và Q8, đồng thời đồng bộ mô tả phương pháp với implementation
thực tế. Cấu hình chính sử dụng bốn expert, routing Top-2, cổng tuyến tính có
context ảnh 6 chiều và backbone MobileNetV3-Small. Các phép so sánh có huấn
luyện dùng năm seed $\{42,43,44,45,46\}$. Báo cáo chỉ dành cho PlantDoc;
kết quả SLIF, nếu được thực hiện, cần trình bày riêng để tránh trộn tập dữ liệu.

Bảng~\ref{tab:overview} liệt kê tám nhóm thực nghiệm. PASS biểu thị tiêu chí
được đáp ứng; TREND biểu thị chênh lệch trung bình theo hướng kỳ vọng nhưng
còn nhỏ so với mức biến thiên thể hiện qua độ lệch chuẩn.

\begin{table}[H]
\caption{Tổng quan tám nhóm thực nghiệm}
\label{tab:overview}
\centering
\footnotesize
\begin{tabular}{@{}clcc@{}}
\toprule
STT & Nhóm thực nghiệm & Train & Trạng thái \\
\midrule
1 & Audit kiến trúc/tham số & Không & \textbf{PASS} \\
2 & Q7: mức sử dụng expert & Không & \textbf{PASS} \\
3 & Q8: shuffled routing & Không & \textbf{PASS} \\
4 & Q1: sparse runtime & Không & \textbf{PASS} \\
5 & Q4: controlled Top-$k$ & Có & \textbf{PASS} \\
6 & Q8: static-uniform control & Có & \textbf{TREND} \\
7 & Ablation context 6-D & Có & \textbf{TREND} \\
8 & Matched MobileNetV3 baseline & Có & \textbf{PASS} \\
\bottomrule
\end{tabular}
\end{table}

\section{Thiết lập thực nghiệm}

\subsection{Dữ liệu và đánh giá}

Tập dữ liệu mặc định là PlantDoc tomato subset. Phân hoạch dữ liệu, số lớp,
số mẫu train/validation/test và quy trình tiền xử lý phải lấy từ log của lần
chạy chính thức; không suy ra từ checkpoint hoặc từ một phiên bản dữ liệu
khác. Các chỉ số chính cho bài toán phân loại là Accuracy và Macro-F1.

\begin{table}[H]
\caption{Thông tin dữ liệu và môi trường chạy}
\label{tab:setup}
\centering
\footnotesize
\begin{tabular}{@{}p{0.44\columnwidth}p{0.47\columnwidth}@{}}
\toprule
Thuộc tính & Giá trị \\
\midrule
Tập dữ liệu & PlantDoc tomato subset \\
Số lớp & 8 \\
Số mẫu train & 2,274 \\
Số mẫu validation & 284 \\
Số mẫu test & 285 \\
Input resolution & $224\times224$ \\
GPU cho huấn luyện & NVIDIA GB10 (CUDA) \\
Môi trường benchmark Q1 & Raspberry Pi 5, CPU aarch64 \\
PyTorch benchmark Q1 & 2.14.0+cpu; không có CUDA \\
Seed & 42, 43, 44, 45, 46 \\
\bottomrule
\end{tabular}
\end{table}

\subsection{Cấu hình mô hình chính}

Mô hình chính có $E=4$ expert với cấu trúc mỗi expert
$576\rightarrow1024\rightarrow576$, sử dụng $k=2$ expert cho mỗi mẫu.
Cụ thể, bộ định tuyến chấm điểm bốn expert và chọn hai expert có điểm cao
nhất để xử lý mẫu; hai expert này được gọi là \emph{tập expert được chọn}.
Cổng routing là cổng tuyến tính context-aware với vector context 6 chiều.
Các phép audit không huấn luyện và các đối chứng phải dùng đúng checkpoint
thuộc nhánh \texttt{moe\_linearcontextaware\_temp0.5}; không thay bằng
checkpoint từ MLP gate.

\subsection{Quy ước báo cáo}

Mọi kết quả phân loại từ các thí nghiệm năm seed được báo dưới dạng
\meanstd{}. Phần phân tích và kết luận chỉ dựa trên trung bình, độ lệch chuẩn
và chênh lệch giữa các giá trị tổng hợp; không phân tích riêng từng seed.
Mọi giá trị trong các bảng dưới đây phải được chép từ đúng tệp CSV được nêu
ở cột ``Nguồn''.

\section{Kết quả}

\subsection{Run 1: Audit kiến trúc và số tham số}

\textbf{Mục đích.} Run 1 là bước kiểm tra tính nhất quán bắt buộc trước khi
thực hiện các phân tích còn lại. Nó xác nhận checkpoint chính thực sự thuộc
mô hình MoE được mô tả trong báo cáo, tránh trường hợp dùng nhầm checkpoint,
nhầm loại gating hoặc nhầm cấu hình số expert và Top-$k$.

\textbf{Cách thực hiện.} Mô hình được instantiate từ code hiện tại rồi nạp
checkpoint đại diện, không huấn luyện lại. Audit kiểm tra metadata
$E$, $k$, kích thước context, cấu trúc từng expert, loại gating, tổng tham số
lưu trữ, số tham số được kích hoạt trong forward pass và FLOPs. Tổng tham số
trainable phải được lấy từ mô hình đã instantiate bằng
$\sum_{p\in\mathcal{P}}\mathrm{numel}(p)$, không dùng tổng tensor trong
\texttt{state\_dict} làm đại diện trực tiếp vì giá trị đó có thể chứa buffer.
Chỉ khi audit đạt yêu cầu mới tiếp tục diễn giải các run phụ thuộc checkpoint.

\begin{table}[H]
\caption{Kết quả audit mô hình}
\label{tab:audit}
\centering
\footnotesize
\begin{tabular}{@{}p{0.48\columnwidth}p{0.24\columnwidth}p{0.17\columnwidth}@{}}
\toprule
Hạng mục & Giá trị & Kiểm tra \\
\midrule
Số expert & 4 & $=4$ \\
Top-$k$ & 2 & $=2$ \\
Kích thước context & 6 & $=6$ \\
Cấu trúc expert & $576$--$1024$--$576$ & $576$--$1024$--$576$ \\
Loại gating & Linear context-aware & Linear context-aware \\
Stored parameters & $5.8535$ M & Xác nhận \\
Activated parameters (Top-2) & $3.4845$ M & Xác nhận \\
FLOPs & $0.0634$ G & Xác nhận \\
Kết luận audit & \textbf{PASS} & Pass \\
\bottomrule
\end{tabular}
\vspace{2pt}

\raggedright\scriptsize Nguồn: \texttt{diagnostics/revision/model\_audit.csv}
và \texttt{diagnostics/revision/active\_complexity\_thop.csv}.
\end{table}

\textbf{Kết luận Run 1: PASS.} Checkpoint phù hợp với cấu hình $E=4$, $k=2$,
\texttt{context\_dim=6}, cấu trúc expert $576$--$1024$--$576$ và cổng
Linear context-aware. Mô hình có $3.4845$ triệu tham số activated và
$5.8535$ triệu tham số lưu trữ; forward pass Top-2 có $0.0634$ GFLOPs.

\subsection{Run 2: Q7--Mức sử dụng expert}

\textbf{Mục đích.} Run 2 kiểm tra bộ định tuyến có phân phối dữ liệu đến cả
bốn expert hay không, qua đó phát hiện hiện tượng \emph{global expert
collapse}, tức một hoặc nhiều expert gần như không bao giờ được chọn. Run này
cũng mô tả sự khác biệt về tần suất kích hoạt expert giữa các lớp bệnh, nhưng
không nhằm chứng minh chuyên môn hóa ngữ nghĩa.

\textbf{Cách thực hiện.} Toàn bộ dữ liệu đánh giá được forward qua năm
checkpoint tương ứng với các seed 42--46, không huấn luyện lại. Với mỗi mẫu,
audit ghi lại các
expert nằm trong Top-2 rồi tổng hợp hai đại lượng. \emph{Global selection
share} là tỷ trọng lượt chọn của mỗi expert trên toàn bộ dữ liệu và có tổng
qua bốn expert bằng 1. \emph{Class activation rate} là xác suất một expert
được kích hoạt trong từng lớp; do mỗi mẫu chọn hai expert, tổng mỗi hàng lớp
bằng $k=2$. Đại lượng thứ hai không được gọi là routing weight.

\begin{table}[H]
\caption{Tỷ lệ lựa chọn chuyên gia toàn cục trên 5 seed}
\label{tab:usage-global}
\centering
\footnotesize
\begin{tabular}{@{}lc@{}}
\toprule
Chuyên gia & Global selection share \\
\midrule
Expert 1 & $0.2446 \pm 0.0237$ \\
Expert 2 & $0.2554 \pm 0.0190$ \\
Expert 3 & $0.2495 \pm 0.0165$ \\
Expert 4 & $0.2505 \pm 0.0134$ \\
Tổng & $1.0000$ \\
\bottomrule
\end{tabular}
\vspace{2pt}

\raggedright\scriptsize Nguồn: \texttt{global\_selection\_share\_summary.csv}.
\end{table}

\begin{table}[H]
\caption{Tỷ lệ kích hoạt theo lớp trên 5 seed}
\label{tab:usage-class}
\centering
\scriptsize
\resizebox{\columnwidth}{!}{%
\begin{tabular}{@{}lcccc@{}}
\toprule
Lớp bệnh & E1 & E2 & E3 & E4 \\
\midrule
Early blight & $0.4600\pm0.4037$ & $0.4100\pm0.3765$ & $0.3300\pm0.3667$ & $0.8000\pm0.2264$ \\
Septoria leaf spot & $0.3524\pm0.4424$ & $0.4095\pm0.2092$ & $0.5048\pm0.3160$ & $0.7333\pm0.3331$ \\
Tomato leaf & $0.4450\pm0.2546$ & $0.5100\pm0.3955$ & $0.4200\pm0.3781$ & $0.6250\pm0.2867$ \\
Bacterial spot & $0.4000\pm0.3092$ & $0.5481\pm0.1420$ & $0.3185\pm0.1563$ & $0.7333\pm0.1841$ \\
Late blight & $0.4000\pm0.4058$ & $0.6273\pm0.3585$ & $0.4455\pm0.3614$ & $0.5273\pm0.2813$ \\
Mosaic virus & $0.6846\pm0.3936$ & $0.5769\pm0.2747$ & $0.5077\pm0.3368$ & $0.2308\pm0.1586$ \\
Yellow virus & $0.5468\pm0.4033$ & $0.4810\pm0.2688$ & $0.6835\pm0.3940$ & $0.2886\pm0.2578$ \\
Mold leaf & $0.5862\pm0.1443$ & $0.6276\pm0.1895$ & $0.4138\pm0.1857$ & $0.3724\pm0.2031$ \\
\bottomrule
\end{tabular}}
\vspace{2pt}

\raggedright\scriptsize Mỗi ô là \meanstd{} tổng hợp từ năm seed. Tổng tỷ lệ
kích hoạt của mỗi hàng bằng $2.0000$, phù hợp với Top-2. Nguồn:
\texttt{class\_activation\_rate\_summary.csv}.
\end{table}

\textbf{Kết luận Run 2: PASS.} Global selection share trung bình của bốn
chuyên gia nằm trong khoảng $0.2446$--$0.2554$, với độ lệch chuẩn
$0.0134$--$0.0237$. Vì tất cả giá trị đều xấp xỉ $0.2500$ và không chuyên gia
nào có mức sử dụng trung bình gần $0.0000$, không có dấu hiệu global expert collapse.
Tuy nhiên, độ lệch chuẩn theo lớp còn khá lớn, nên kết quả tổng hợp này chưa
đủ để khẳng định mỗi chuyên gia đã học một vai trò ngữ nghĩa bệnh ổn định.

\subsection{Run 3: Q8--Shuffled-routing counterfactual}

\textbf{Mục đích.} Run 3 kiểm tra liệu ánh xạ phụ thuộc đầu vào
$\text{input}\rightarrow\text{tập expert được chọn}$ có thực sự đóng góp vào
hiệu năng hay không. Đây là đối chứng phản thực tế: nếu chỉ cần có nhiều
expert mà không cần lựa chọn đúng expert cho từng mẫu, việc phá vỡ ánh xạ này
sẽ không làm giảm hiệu năng trung bình.

\textbf{Cách thực hiện.} Thí nghiệm giữ nguyên kiến trúc, trọng số và đầu ra
của bộ định tuyến, nhưng xáo trộn quyết định định tuyến giữa các input; nói
cách khác, một mẫu không còn được xử lý bởi tập expert Top-2 vốn được chọn cho
chính mẫu đó. Mỗi checkpoint được thực hiện 100 lần xáo trộn, không huấn
luyện lại. Kết quả cuối cùng của learned routing và shuffled routing đều được
tổng hợp dưới dạng trung bình $\pm$ độ lệch chuẩn từ năm checkpoint. Hai độ
chênh chính là
\begin{align}
\Delta_{\mathrm{Acc}} &= \mathrm{Acc}_{\mathrm{learned}}
- \mathrm{Acc}_{\mathrm{shuffled}},\\
\Delta_{\mathrm{F1}} &= \mathrm{F1}_{\mathrm{learned}}
- \mathrm{F1}_{\mathrm{shuffled}}.
\end{align}

\begin{table}[H]
\caption{Tổng hợp shuffled routing trên năm seed}
\label{tab:shuffle-summary}
\centering
\scriptsize
\begin{tabular}{@{}lccc@{}}
\toprule
Metric & Learned & Shuffled & Chênh lệch \\
\midrule
Accuracy & $0.8596\pm0.0124$ & $0.8287\pm0.0102$ & $+0.0309$ \\
Macro-F1 & $0.8275\pm0.0153$ & $0.7912\pm0.0129$ & $+0.0363$ \\
\bottomrule
\end{tabular}
\vspace{2pt}

\raggedright\scriptsize Learned và Shuffled được báo dưới dạng \meanstd{}
tổng hợp từ năm seed; mỗi checkpoint dùng 100 lần hoán vị. Nguồn:
\texttt{shuffle\_seed\_summary.csv}.
\end{table}

\textbf{Kết luận Run 3: PASS.} Khi phá vỡ quan hệ giữa ảnh đầu vào và tập
expert được chọn, Accuracy trung bình giảm từ $0.8596\pm0.0124$ xuống
$0.8287\pm0.0102$, còn Macro-F1 giảm từ $0.8275\pm0.0153$ xuống
$0.7912\pm0.0129$. Các chênh lệch trung bình tương ứng là $0.0309$ và
$0.0363$, cho thấy định tuyến phụ thuộc đầu vào có đóng góp vào hiệu năng mô
hình. Kết quả này không phải bằng chứng rằng từng chuyên gia đã học được một
ý nghĩa bệnh riêng biệt.

\subsection{Run 4: Q1--Dynamic sparse runtime}

\textbf{Mục đích.} Run 4 đo quan hệ giữa số expert được kích hoạt và thời gian
thực thi của mô hình bằng PyTorch eager. Khi Top-$k$ tăng, số expert tham gia
forward pass tăng và thời gian thực thi được kỳ vọng tăng tương ứng. Run này
không đánh giá lại chất lượng phân loại và không huấn luyện mô hình.

\textbf{Cách thực hiện.} Cùng một checkpoint và cùng bộ trọng số được dùng
cho $k\in\{1,2,3,4\}$; script chỉ thay đổi số expert active nên kích thước và
số tham số lưu trữ không đổi. Benchmark được thực hiện trực tiếp bằng
PyTorch eager trên Raspberry Pi 5 với CPU aarch64, PyTorch 2.14.0+cpu và
không có CUDA; mô hình không được xuất sang ONNX cho phép đo này. Mỗi cấu
hình dùng batch size 1, một CPU thread, 30 lượt warm-up và 300 lượt đo. Input
tensor được chuẩn bị sẵn; thời gian chỉ tính forward pass, không gồm đọc ảnh,
tiền xử lý hoặc I/O. Các đại lượng được báo gồm mean, median, P95 và độ lệch
chuẩn.

\begin{table}[H]
\caption{PyTorch eager runtime trên Raspberry Pi 5}
\label{tab:runtime}
\centering
\footnotesize
\begin{tabular}{@{}crrrr@{}}
\toprule
Top-$k$ & Mean (ms) & Median (ms) & P95 (ms) & Std. (ms) \\
\midrule
1 & $23.8236$ & $23.7939$ & $24.0175$ & $0.1520$ \\
\textbf{2} & \textbf{25.4065} & \textbf{25.3656} & \textbf{25.6108} & \textbf{0.3042} \\
3 & $27.4678$ & $27.3820$ & $27.6219$ & $1.0129$ \\
4 & $31.6735$ & $31.6358$ & $31.9404$ & $0.1958$ \\
\bottomrule
\end{tabular}
\vspace{2pt}

\raggedright\scriptsize Batch size 1, một CPU thread, 30 lượt warm-up và 300
lượt đo cho mỗi $k$. Thời gian chỉ gồm forward pass. Nguồn:
\texttt{pi\_runtime/sparse\_runtime\_pi.csv}.
\end{table}

\textbf{Kết luận Run 4: PASS.} Latency trung bình tăng đơn điệu từ
$23.8236$ ms tại $k=1$ lên $31.6735$ ms tại $k=4$, tương ứng tăng
$32.9499\%$. Cấu hình chính Top-2 đạt $25.4065$ ms, cao hơn Top-1
$6.6440\%$. Kết quả cho thấy trong phép đo bằng PyTorch eager, thời gian thực
thi tăng khi số expert được kích hoạt tăng. Kết quả không bao gồm đọc ảnh,
tiền xử lý hoặc I/O và không đại diện cho ONNX runtime.

\subsection{Run 5: Q4--Controlled Top-\texorpdfstring{$k$}{k} ablation}

\textbf{Mục đích.} Run 5 xác định Top-2 có phải operating point hợp lý giữa
độ chính xác và chi phí tính toán hay không. Khác với Run 4 chỉ thay đổi $k$
khi benchmark cùng trọng số, Run này huấn luyện riêng từng cấu hình để đo
đúng ảnh hưởng của $k$ đến chất lượng mô hình sau tối ưu hóa.

\textbf{Cách thực hiện.} Thí nghiệm giữ cố định $E=4$, cổng Linear
context-aware, backbone, expert, dữ liệu và thiết lập huấn luyện; chỉ thay đổi
$k=1,2,3,4$. Mỗi giá trị $k$ được huấn luyện với năm seed 42--46, tạo tổng
cộng 20 run. Accuracy và Macro-F1 được báo bằng trung bình $\pm$ độ lệch
chuẩn của năm seed rồi đối chiếu với runtime của Run 4 để phân tích
trade-off. Bảng~\ref{tab:topk} phải lấy từ kết
quả revision mới, không dùng tệp kết quả cũ thuộc gate hoặc cấu hình khác.

\begin{table}[H]
\caption{Ablation Top-$k$ có kiểm soát}
\label{tab:topk}
\centering
\footnotesize
\begin{tabular}{@{}ccc@{}}
\toprule
$k$ & Accuracy & Macro-F1 \\
\midrule
1 & $0.8309\pm0.0169$ & $0.7927\pm0.0126$ \\
2 & \best{$0.8596\pm0.0124$} & \best{$0.8275\pm0.0153$} \\
3 & $0.8540\pm0.0231$ & $0.8243\pm0.0268$ \\
4 & $0.8491\pm0.0197$ & $0.8124\pm0.0291$ \\
\bottomrule
\end{tabular}
\vspace{2pt}

\raggedright\scriptsize Accuracy và Macro-F1 được báo dưới dạng \meanstd{}
qua 5 seed. Nguồn: \texttt{diagnostics/revision/topk/topk\_summary.csv}.
\end{table}

\textbf{Operating point được chọn:} $k=2$. \textbf{Kết luận Run 5: PASS.}
Top-2 đạt Accuracy $0.8596\pm0.0124$ và Macro-F1 $0.8275\pm0.0153$. Top-3
cho kết quả trung bình gần Top-2 nhưng dùng thêm một expert và có độ lệch
chuẩn lớn hơn; Top-4 không cải thiện hiệu năng trung bình dù kích hoạt toàn
bộ expert. Vì vậy,
Top-2 là operating point hợp lý trong các cấu hình đã thử.

\subsection{Run 6: Q8--Static-uniform 4-expert control}

\textbf{Mục đích.} Run 6 tách đóng góp của routing phụ thuộc đầu vào khỏi lợi
ích đơn thuần do mô hình lưu trữ bốn expert. Đây là đối chứng matched-capacity
trực tiếp cho Q8: nếu Main MoE tốt hơn ổn định, cải thiện khó có thể chỉ được
giải thích bằng số tham số expert được lưu trữ.

\textbf{Cách thực hiện.} Đối chứng static-uniform giữ cùng backbone, bốn
expert $576\rightarrow1024\rightarrow576$, residual path và classifier cuối,
nhưng bỏ input-dependent gate và kết hợp expert bằng trọng số cố định đồng
đều. Năm checkpoint control được huấn luyện độc lập với các seed 42--46.
Sau huấn luyện, Accuracy và Macro-F1 của mỗi cấu hình được tổng hợp thành
trung bình $\pm$ độ lệch chuẩn trên năm seed trong
\texttt{control\_summary.csv}.

\textbf{Cách diễn giải.} So sánh trung bình và độ lệch chuẩn cho biết learned
routing có cải thiện hiệu năng tổng hợp so với kết hợp expert bằng trọng số
cố định hay không, trong điều kiện hai cấu hình có dung lượng lưu trữ gần
tương đương.

\begin{table}[H]
\caption{So sánh Main MoE với static-uniform control}
\label{tab:static-summary}
\centering
\footnotesize
\begin{tabular}{@{}lcc@{}}
\toprule
Mô hình & Accuracy & Macro-F1 \\
\midrule
Main MoE & \textbf{$0.8596\pm0.0124$} & \textbf{$0.8275\pm0.0153$} \\
Static-uniform 4-expert & $0.8365\pm0.0234$ & $0.8019\pm0.0247$ \\
\textbf{MoE $-$ Static} & \textbf{$+0.0232$} & \textbf{$+0.0256$} \\
\bottomrule
\end{tabular}
\vspace{2pt}

\raggedright\scriptsize Kết quả là \meanstd{} tổng hợp từ năm seed. Nguồn:
\texttt{control\_summary.csv}.
\end{table}

\textbf{Kết luận Run 6: TREND.} Main MoE đạt Accuracy
$0.8596\pm0.0124$ và Macro-F1 $0.8275\pm0.0153$, so với
$0.8365\pm0.0234$ và $0.8019\pm0.0247$ của static-uniform. Chênh lệch trung
bình lần lượt là $0.0232$ và $0.0256$, đồng thời độ lệch chuẩn của Main MoE
thấp hơn trên cả hai metric. Vì hai mô hình cùng sử dụng
backbone và bốn expert có cấu trúc tương đương, các kết quả tổng hợp này cho
thấy routing phụ thuộc đầu vào cải thiện hiệu năng trung bình và giảm biến
thiên so với static-uniform.

\subsection{Run 7: Ablation context ảnh 6 chiều}

\textbf{Mục đích.} Run 7 định lượng mức đóng góp của sáu đặc trưng context
suy ra từ ảnh được đưa vào cổng routing. Ablation này đặc biệt cần thiết vì
implementation chính sử dụng nhánh context 6-D trong khi mô tả cũ của
manuscript chưa phản ánh đầy đủ nhánh này.

\textbf{Cách thực hiện.} Mô hình đối chứng giữ nguyên backbone, kiến trúc MoE,
số expert, Top-2 và số tham số, nhưng đặt toàn bộ vector context về 0 trong cả
train, validation và test. Zero-context MoE được huấn luyện trên năm seed
42--46. Accuracy và Macro-F1 của hai cấu hình được tổng hợp thành trung bình
$\pm$ độ lệch chuẩn để đánh giá ảnh hưởng của tín hiệu context. Nhánh context
6 chiều vẫn phải được mô tả đúng vì checkpoint chính sử dụng cổng
context-aware.

\begin{table}[H]
\caption{Ảnh hưởng của context 6 chiều}
\label{tab:context}
\centering
\footnotesize
\begin{tabular}{@{}lcc@{}}
\toprule
Cấu hình & Accuracy & Macro-F1 \\
\midrule
Main MoE (context 6-D) & \best{$0.8596\pm0.0124$} & \best{$0.8275\pm0.0153$} \\
Zero-context MoE & $0.8330\pm0.0343$ & $0.8018\pm0.0342$ \\
\bottomrule
\end{tabular}
\vspace{2pt}

\raggedright\scriptsize Các ô hiệu năng báo \meanstd{} qua 5 seed. Nguồn:
\texttt{context\_seed\_metrics.csv}.
\end{table}

\textbf{Kết luận Run 7:} Main MoE sử dụng context 6-D đạt Accuracy
$0.8596\pm0.0124$ và Macro-F1 $0.8275\pm0.0153$, cao hơn zero-context lần
lượt $0.0267$ và $0.0257$. Độ lệch chuẩn của context-aware cũng thấp hơn trên
cả Accuracy và Macro-F1, cho thấy độ biến thiên tổng hợp thấp hơn. Kết quả
này cho thấy sáu đặc trưng context có đóng góp tích cực vào hiệu năng trung
bình của mô hình. Vì checkpoint chính sử dụng cổng context-aware, nhánh
context 6 chiều phải được mô tả đầy đủ trong manuscript.

\subsection{Run 8: Matched MobileNetV3-Small baseline}

\begin{table}[H]
\caption{Kết quả MobileNetV3-Small pretrained cùng họ torchvision}
\label{tab:baseline}
\centering
\footnotesize
\begin{tabular}{@{}lcc@{}}
\toprule
Mô hình & Accuracy & Macro-F1 \\
\midrule
MobileNetV3-Small & $0.8267\pm0.0223$ & $0.7894\pm0.0242$ \\
\bottomrule
\end{tabular}
\vspace{2pt}

\raggedright\scriptsize Nguồn: \texttt{diagnostics/revision/controls/control\_summary.csv}.
\end{table}

\section{Thảo luận tổng hợp}

Các giá trị tổng hợp cho thấy Main MoE có Accuracy và Macro-F1 trung bình cao
hơn static-uniform 4-expert, đồng thời learned routing cao hơn shuffled
routing trên cả hai metric. Kết hợp hai kết quả này cho thấy routing phụ
thuộc đầu vào đóng góp vượt ra ngoài lợi ích do dung lượng expert được lưu
trữ. Kết quả này vẫn không phải bằng chứng trực tiếp rằng các expert đã học
được chuyên môn hóa ngữ nghĩa.

\begin{table}[H]
\caption{Ma trận bằng chứng và kết luận cuối}
\label{tab:evidence}
\centering
\footnotesize
\begin{tabular}{@{}p{0.53\columnwidth}p{0.35\columnwidth}@{}}
\toprule
Điều kiện & Kết quả \\
\midrule
Audit checkpoint/architecture đạt yêu cầu & \textbf{PASS} \\
Runtime tăng hợp lý theo $k$ & \textbf{PASS} \\
Top-2 là operating point hợp lý & \textbf{PASS} \\
Không có expert collapse toàn cục & \textbf{PASS} \\
Learned routing $>$ shuffled routing & \textbf{PASS} \\
Main MoE $>$ static-uniform control & \textbf{TREND} \\
Context 6-D tạo cải thiện trung bình & \textbf{TREND} \\
\midrule
Kết luận tổng hợp & \textbf{PASS có giới hạn} \\
\bottomrule
\end{tabular}
\end{table}

\section{Hạn chế}

Các giá trị trung bình và độ lệch chuẩn chỉ được ước lượng từ năm seed nên độ
chắc chắn của phân phối thực nghiệm còn hạn chế. Benchmark latency phụ thuộc
phần cứng, power mode, số thread và phiên bản PyTorch; do đó chỉ so sánh các
cấu hình được đo trong cùng điều kiện. Phép đo runtime trong báo cáo chỉ đại
diện cho PyTorch eager trên Raspberry Pi 5, không phải ONNX runtime. Các kết
luận của báo cáo hiện giới hạn ở PlantDoc và không tự động khái quát sang SLIF
hoặc tập dữ liệu khác.

\section{Kết luận}

Báo cáo đã tổng hợp tám nhóm thực nghiệm revision và gắn từng bảng với đầu ra
CSV tương ứng. Top-2 là operating point tốt nhất trong ablation có kiểm soát,
đồng thời learned routing tốt hơn shuffled routing và không xuất hiện global
expert collapse. Nhánh context 6-D cho Accuracy và Macro-F1 trung bình cao
hơn zero-context, đồng thời có độ lệch chuẩn thấp hơn trên cả hai metric. Các
kết luận hiện giới hạn ở PlantDoc và cần giữ nguyên các lưu ý về số lượng
seed, runtime và implementation.

\end{document}
